% SOURCE_STATEMENT: "The fact that RMT, and the Tracy-Widom distributions, arise in so many problems in so many different areas leads one to the following question: how can one characterize RMT in purely probabilistic terms? For example, we know that if we take i.i.d.'s $(a_1,a_2,...)$, add them up, and then center and scale appropriately, $$ (a_1,a_2,...)\\to(S_1,S_2,...),\\qquad S_n=\\frac{\\sum_{i=1}^n a_i-n\\mu}{\\sqrt n}, $$ then as $n\\to\\infty$, $S_n$ converges in distribution to a Gaussian random variable: this is the famous Central Limit Theorem. The analogous situation for RMT is the following: take i.i.d.'s $(a_1,a_2,...)$, perform an operation $X$ on them, $$ (a_1,a_2,...)\\to(X_1,X_2,...), $$ and as $n\\to\\infty$ the $X_n$'s converge to the Gaudin distribution, or the Tracy-Widom distribution. The question is, ``What is $X$?'' Important progress towards answering this question has been made recently, and independently, by Baik-Suidan and Bodineau-Martin, but the full problem remains open and very challenging."
% Preparation and final self-audit: CONTRIBUTING.md prompts C and D; source/audit artifacts retained in the local physics_statements_2026_10_11 batch.
% !TeX program = lualatex
\documentclass[11pt]{article}
\usepackage[margin=25mm]{geometry}
\usepackage{fontspec}
\setmainfont{Latin Modern Roman}
\usepackage{amsmath,amssymb,mathtools}
\usepackage{unicode-math}
\setmathfont{Latin Modern Math}
\usepackage{xurl,hyperref}
\hypersetup{colorlinks=true,urlcolor=blue}
\setcounter{secnumdepth}{0}
\setlength{\parindent}{0pt}
\setlength{\parskip}{5pt}
\setlength{\emergencystretch}{3em}
\newcommand{\sref}[2]{\hyperlink{source:#1:#2}{[S#2]}}
\newcommand{\cataloguescope}[1]{\paragraph{Recorded target.}#1}
\newcommand{\R}{\mathbb R}
\newcommand{\C}{\mathbb C}
\newcommand{\N}{\mathbb N}
\newcommand{\Z}{\mathbb Z}
\newcommand{\Q}{\mathbb Q}
\newcommand{\F}{\mathbb F}
\newcommand{\eps}{\varepsilon}
\begin{document}
\section{A Tracy-Widom Central Limit Theorem}\label{statement}

\subsection{Definitions and mathematical statement}

\paragraph{Shared notation.}$\mathbb N=\{1,2,\ldots\}$; the standard number fields have their usual meanings. All additional conventions are stated locally.


\paragraph{Definitions and notation.}Let $(a_i)_{i\ge1}$ be independent identically distributed random variables. The ordinary central limit theorem studies centered, rescaled sums. Deift instead asks for probabilistically natural operations yielding random-matrix limits. The Gaudin law describes a random-matrix spacing distribution; Tracy--Widom laws describe edge fluctuations in the corresponding symmetry classes. Finding a useful intrinsic operation includes specifying its inputs, scaling and admissible laws.

\paragraph{Catalogue wording retained for traceability.}
Characterize random-matrix statistics in purely probabilistic terms: identify a natural operation $X$ on i.i.d. input $(a_1,a_2,\ldots)$ producing normalized variables $(X_n)$ whose distributions tend to the Gaudin or Tracy--Widom law. Determine conditions and normalization giving a meaningful universality principle, rather than just forcing the desired law by an arbitrary transformation.
\paragraph{Scope and source qualification.}The original asks What is X rather than fixing last-passage percolation as the only operation. The Gaudin target as well as Tracy–Widom is retained. It is a characterization programme, not a statement that every i.i.d. law has such an edge limit without hypotheses.
\subsection{Short English statement}
Find a probabilistic analogue of the central limit theorem whose limit is a random-matrix spacing or edge law.
\subsection{Sources}
\begin{itemize}
\item[\hypertarget{source:statement:1}{S1}] ULAM.AI and the UnsolvedMath Contributors, UnsolvedMath, catalogue entry. Catalogue curation under CC BY 4.0; definitions and English exposition added locally.
\url{https://huggingface.co/datasets/ulamai/UnsolvedMath}.
\item[\hypertarget{source:statement:2}{S2}] Percy Deift, Some open problems in random matrix theory and the theory of integrable systems (2007), arXiv:0712.0849, Problem 10.
\url{https://arxiv.org/abs/0712.0849}.
\end{itemize}
\label{end:statement}
\end{document}
